\par\noindent\textbf{(1)} 总自旋角动量近似为 \(\boldsymbol L=S\boldsymbol n\)。力矩由取向势变分给 \[\boldsymbol\tau=-\boldsymbol n\times\partial U/\partial\boldsymbol n=\kappa(\boldsymbol n\cdot\boldsymbol k)(\boldsymbol n\times\boldsymbol k)\]。故 \(S\dot{\boldsymbol n}=\kappa\cos\theta(\boldsymbol n\times\boldsymbol k)\)，法向长度 \(|\boldsymbol n|=1\) 自动保持。
\par\noindent\textbf{(2)} 因为 \(\dot{\boldsymbol n}\perp\boldsymbol k\)，\(\boldsymbol n\cdot\boldsymbol k=\cos\theta\) 守恒。若 \(\dot{\boldsymbol n}=\Omega_p(\boldsymbol k\times\boldsymbol n)\)，与上式比较得 \(\Omega_p=-\kappa\cos\theta/S\)；进动方向取决于符号约定和倾角，近赤道处速率趋零。
\par\noindent\textbf{(3)} 瞬时功率为 \(P=\boldsymbol\tau\cdot\boldsymbol\omega\)。对势随时间的变化，\(\dot U=-\kappa(\boldsymbol n\cdot\boldsymbol k)(\dot{\boldsymbol n}\cdot\boldsymbol k)=0\)，说明理想轴对称力矩只重定向自旋，不改变自旋能量；不能把进动与摩擦损失混为一谈。
\par\noindent\textbf{(4)} 若保留横向惯量 \(I_\perp\)，角动力学包含 \(I_\perp\ddot{\boldsymbol n}\) 的快速项，除低频进动外出现尺度约为 \(\omega_{\mathrm{nut}}\sim S/I_\perp\) 的章动。快自旋近似的适用条件是 \(\kappa I_\perp/S^2\ll1\)，在失效时必须使用完整Euler角动力学。
